% !TEX root = ../main.tex
% =====================================================================
% APPENDIX A -- ADDITIONAL RESULTS          target: 3 - 5 pages
%
% Everything that supports the main story but is not needed to follow it.
% The rule I apply: if removing it would break an argument in section 8,
% it belongs in section 8; otherwise it belongs here.
%
% These pages COUNT towards the 20-30 page limit, so keep the appendix
% tight: tables and figures with one or two sentences each, no long prose.
% =====================================================================

\section{Additional Results}
\label{app:additional}

\begin{guide}
Structure of this appendix:
\begin{enumerate}\itemsep2pt
  \item full leaderboard (every model I trained),
  \item extended ablations,
  \item per-dataset results,
  \item extra experiments (ensembles, hyperparameter search, discarded approaches),
  \item extra figures.
\end{enumerate}
Every item needs one sentence saying why it is here. An appendix with no text at all
looks like a data dump.
\end{guide}

% ---------------------------------------------------------------------
\subsection{Full leaderboard}
\label{app:leaderboard}

{Every encoder $\times$ pooling combination I trained, ranked by \webtraffic{}
accuracy; \Cref{tab:web_main} in the main text keeps only the two headline variants
and the baselines.}

% \begin{table}[t]
\centering
\small
\caption{Complete leaderboard: all 72 evaluated models, ranked by WebTraffic accuracy. The Model column is the configuration identifier under configs/models/, so any model named in the report can be looked up here. Encoder and pooling names are prefixed sea\_ for components introduced in this work and mil\_ for components reused unchanged from MILLET. A dash marks a model screened on WebTraffic only.}
\label{tab:full_leaderboard}
\begin{tabular}{lrrrrrrrrrrrr}
\toprule
\# & Model & Encoder & Pooling & Origin & Acc. & AOPCR & NDCG & Loss & \#DS & UCR-85 Acc. & Params & Size (MB) \\
\midrule
1 & seanet\_gated\_mean\_topk & sea\_mstcn\_sep\_gated & sea\_topk\_conjunctive & ours & 0.955 & 2.23 & 0.750 & 0.254 & 85 & 0.8083 & 61,740 & 1.50 \\
2 & seanet\_conjunctive & sea\_mstcn\_sep & mil\_conjunctive & half-ours & 0.954 & 1.50 & 0.698 & 0.262 & 85 & 0.8238 & 269,083 & 3.54 \\
3 & seanet\_gated\_max\_topk & sea\_mstcn\_sep\_gated & sea\_topk\_conjunctive & ours & 0.952 & 2.27 & 0.719 & 0.308 & 85 & 0.8108 & 61,740 & 1.50 \\
4 & seanet\_topk\_nofocus & sea\_mstcn\_sep\_bottleneck & sea\_topk\_conjunctive & ours & 0.950 & 2.30 & 0.765 & 0.263 & 0 & -- & 41,324 & 1.43 \\
5 & seanet\_spiketrend\_topk & sea\_mstcn\_sep\_spiketrend & sea\_topk\_conjunctive & ours & 0.950 & 2.32 & 0.756 & 0.288 & 85 & 0.8089 & 67,020 & 1.52 \\
6 & seanet\_inputgate\_mschan & sea\_multiscale\_channels & sea\_topk\_conjunctive & ours & 0.948 & 2.32 & 0.757 & 0.318 & 0 & -- & 69,768 & 1.54 \\
7 & seanet\_gated\_mschan & sea\_multiscale\_channels & sea\_topk\_conjunctive & ours & 0.948 & 2.22 & 0.717 & 0.289 & 0 & -- & 67,592 & 1.53 \\
8 & seanet\_bottleneck\_topk & sea\_mstcn\_sep\_bottleneck & sea\_topk\_conjunctive & ours & 0.947 & 2.62 & 0.772 & 0.282 & 85 & 0.8097 & 41,324 & 1.43 \\
9 & seanet\_recon\_adaptive & sea\_mstcn\_sep\_recon & sea\_adaptive\_classwise & ours & 0.946 & 2.60 & 0.768 & 0.296 & 85 & 0.8140 & 62,935 & 1.51 \\
10 & seanet\_inputgate\_topk & sea\_mstcn\_sep\_inputgate & sea\_topk\_conjunctive & ours & 0.946 & 2.80 & 0.748 & 0.261 & 85 & 0.8131 & 58,092 & 1.49 \\
11 & seanet\_classwise & sea\_mstcn\_sep & sea\_classwise\_conjunctive & ours & 0.946 & 1.72 & 0.686 & 0.292 & 85 & 0.8292 & 269,164 & 3.54 \\
12 & seanet\_gated\_max & sea\_mstcn\_sep\_gated & sea\_classwise\_conjunctive & ours & 0.944 & 1.89 & 0.592 & 0.306 & 85 & 0.8237 & 61,740 & 1.50 \\
13 & seanet\_gated\_last\_topk & sea\_mstcn\_sep\_gated & sea\_topk\_conjunctive & ours & 0.942 & 2.31 & 0.748 & 0.328 & 85 & 0.8015 & 61,740 & 1.50 \\
14 & seanet & sea\_mstcn\_sep & mil\_additive & half-ours & 0.942 & 1.58 & 0.729 & 0.260 & 85 & 0.8298 & 269,083 & 3.54 \\
15 & seanet\_gated\_last & sea\_mstcn\_sep\_gated & sea\_classwise\_conjunctive & ours & 0.942 & 2.52 & 0.742 & 0.327 & 85 & 0.7888 & 61,740 & 1.50 \\
16 & seanet\_recon\_topk & sea\_mstcn\_sep\_recon & sea\_topk\_conjunctive & ours & 0.940 & 2.51 & 0.698 & 0.321 & 85 & 0.8063 & 62,925 & 1.51 \\
17 & seanet\_topk\_k025 & sea\_mstcn\_sep\_bottleneck & sea\_topk\_conjunctive & ours & 0.940 & 2.65 & 0.733 & 0.277 & 0 & -- & 41,324 & 1.43 \\
18 & seanet\_bottleneck\_softmax & sea\_mstcn\_sep\_bottleneck & sea\_softmax\_conjunctive & ours & 0.940 & 1.68 & 0.740 & 0.345 & 85 & 0.8106 & 41,325 & 1.43 \\
19 & seanet\_gated\_last\_adaptive & sea\_mstcn\_sep\_gated & sea\_adaptive\_classwise & ours & 0.938 & 2.39 & 0.793 & 0.289 & 85 & 0.8029 & 61,750 & 1.50 \\
20 & seanet\_spiketrend\_adaptive & sea\_mstcn\_sep\_spiketrend & sea\_adaptive\_classwise & ours & 0.934 & 1.91 & 0.747 & 0.301 & 85 & 0.8010 & 67,030 & 1.53 \\
21 & seanet\_recon\_softmax & sea\_mstcn\_sep\_recon & sea\_softmax\_conjunctive & ours & 0.932 & 1.64 & 0.723 & 0.341 & 0 & -- & 62,926 & 1.51 \\
22 & seanet\_slim\_topk & sea\_mstcn\_sep & sea\_topk\_conjunctive & ours & 0.932 & 2.69 & 0.778 & 0.306 & 85 & 0.8093 & 57,580 & 1.48 \\
23 & seanet\_bottleneck\_gatedattn & sea\_mstcn\_sep\_bottleneck & sea\_gated\_attention & ours & 0.930 & 2.23 & 0.741 & 0.313 & 85 & 0.8146 & 41,844 & 1.43 \\
24 & seanet\_spiketrend\_gatedattn & sea\_mstcn\_sep\_spiketrend & sea\_gated\_attention & ours & 0.930 & 1.59 & 0.721 & 0.350 & 85 & 0.8039 & 67,540 & 1.53 \\
25 & seanet\_gated & sea\_mstcn\_sep\_gated & sea\_classwise\_conjunctive & ours & 0.930 & 2.41 & 0.637 & 0.359 & 85 & 0.7900 & 61,740 & 1.50 \\
26 & seanet\_slim\_adaptive & sea\_mstcn\_sep & sea\_adaptive\_classwise & ours & 0.930 & 2.45 & 0.793 & 0.336 & 0 & -- & 57,590 & 1.49 \\
27 & seanet\_bottleneck\_mschan & sea\_multiscale\_channels & sea\_topk\_conjunctive & ours & 0.926 & 2.24 & 0.777 & 0.357 & 0 & -- & 47,176 & 1.45 \\
28 & seanet\_bottleneck\_mschan\_lean & sea\_multiscale\_channels & sea\_topk\_conjunctive & ours & 0.924 & 2.76 & 0.735 & 0.389 & 0 & -- & 43,576 & 1.44 \\
29 & seanet\_gated\_last\_softmax & sea\_mstcn\_sep\_gated & sea\_softmax\_conjunctive & ours & 0.920 & 1.99 & 0.738 & 0.381 & 0 & -- & 61,741 & 1.50 \\
30 & millet\_paper & mil\_inceptiontime & mil\_conjunctive & millet & 0.920 & 13.27 & 0.661 & 0.371 & 84 & 0.8434 & 423,707 & 4.11 \\
31 & seanet\_acp & sea\_mstcn\_sep & sea\_adaptive\_classwise & ours & 0.918 & 1.61 & 0.619 & 0.360 & 85 & 0.8134 & 269,174 & 3.54 \\
32 & seanet\_gated\_last\_gatedattn & sea\_mstcn\_sep\_gated & sea\_gated\_attention & ours & 0.918 & 1.35 & 0.757 & 0.365 & 85 & 0.8048 & 62,260 & 1.50 \\
33 & seanet\_slim & sea\_mstcn\_sep & sea\_classwise\_conjunctive & ours & 0.916 & 2.48 & 0.719 & 0.383 & 0 & -- & 57,580 & 1.48 \\
34 & seanet\_gated\_max\_softmax & sea\_mstcn\_sep\_gated & sea\_softmax\_conjunctive & ours & 0.914 & 1.88 & 0.776 & 0.372 & 0 & -- & 61,741 & 1.50 \\
35 & seanet\_inputgate & sea\_mstcn\_sep\_inputgate & sea\_classwise\_conjunctive & ours & 0.914 & 2.48 & 0.715 & 0.398 & 0 & -- & 58,092 & 1.49 \\
36 & seanet\_gated\_mean\_softmax & sea\_mstcn\_sep\_gated & sea\_softmax\_conjunctive & ours & 0.912 & 2.09 & 0.765 & 0.368 & 0 & -- & 61,741 & 1.50 \\
37 & seanet\_recon & sea\_mstcn\_sep\_recon & sea\_classwise\_conjunctive & ours & 0.910 & 2.50 & 0.680 & 0.417 & 0 & -- & 62,925 & 1.51 \\
38 & seanet\_topk\_k100 & sea\_mstcn\_sep\_bottleneck & sea\_topk\_conjunctive & ours & 0.908 & 2.44 & 0.722 & 0.395 & 0 & -- & 41,324 & 1.43 \\
39 & seanet\_gated\_max\_adaptive & sea\_mstcn\_sep\_gated & sea\_adaptive\_classwise & ours & 0.906 & 2.07 & 0.729 & 0.444 & 0 & -- & 61,750 & 1.50 \\
40 & seanet\_inputgate\_adaptive & sea\_mstcn\_sep\_inputgate & sea\_adaptive\_classwise & ours & 0.905 & 2.65 & 0.732 & 0.390 & 85 & 0.8153 & 58,102 & 1.49 \\
41 & seanet\_spiketrend & sea\_mstcn\_sep\_spiketrend & sea\_classwise\_conjunctive & ours & 0.904 & 2.15 & 0.699 & 0.395 & 0 & -- & 67,020 & 1.52 \\
42 & seanet\_bottleneck & sea\_mstcn\_sep\_bottleneck & sea\_classwise\_conjunctive & ours & 0.902 & 2.57 & 0.733 & 0.395 & 0 & -- & 41,324 & 1.43 \\
43 & seanet\_gated\_mean & sea\_mstcn\_sep\_gated & sea\_classwise\_conjunctive & ours & 0.902 & 2.49 & 0.714 & 0.393 & 0 & -- & 61,740 & 1.50 \\
44 & seanet\_inputgate\_softmax & sea\_mstcn\_sep\_inputgate & sea\_softmax\_conjunctive & ours & 0.894 & 1.64 & 0.696 & 0.454 & 0 & -- & 58,093 & 1.49 \\
45 & seanet\_slim\_gatedattn & sea\_mstcn\_sep & sea\_gated\_attention & ours & 0.894 & 2.30 & 0.714 & 0.403 & 0 & -- & 58,100 & 1.49 \\
46 & seanet\_recon\_attnmax & sea\_mstcn\_sep\_recon & sea\_attention\_max & ours & 0.892 & 2.32 & 0.725 & 0.499 & 0 & -- & 62,935 & 1.51 \\
47 & seanet\_slim\_softmax & sea\_mstcn\_sep & sea\_softmax\_conjunctive & ours & 0.890 & 1.83 & 0.705 & 0.419 & 0 & -- & 57,581 & 1.49 \\
48 & seanet\_slim\_gatedattn & sea\_mstcn\_sep\_gated & mil\_attention & half-ours & 0.888 & 2.40 & 0.710 & 0.426 & 0 & -- & 58,100 & 1.49 \\
49 & seanet\_topk\_k005 & sea\_mstcn\_sep\_bottleneck & sea\_topk\_conjunctive & ours & 0.888 & 2.29 & 0.715 & 0.401 & 0 & -- & 41,324 & 1.43 \\
50 & millet & mil\_inceptiontime & mil\_conjunctive & millet & 0.887 & 2.57 & 0.677 & 0.507 & 84 & 0.8274 & 423,707 & 4.11 \\
51 & seanet\_gated\_max\_attnmax & sea\_mstcn\_sep\_gated & sea\_attention\_max & ours & 0.886 & 1.12 & 0.575 & 0.588 & 0 & -- & 61,750 & 1.50 \\
52 & seanet\_topk\_k050 & sea\_mstcn\_sep\_bottleneck & sea\_topk\_conjunctive & ours & 0.882 & 3.11 & 0.732 & 0.418 & 0 & -- & 41,324 & 1.43 \\
53 & seanet\_softmax & sea\_mstcn\_sep & sea\_softmax\_conjunctive & ours & 0.880 & 1.02 & 0.605 & 0.455 & 85 & 0.8131 & 269,165 & 3.54 \\
54 & seanet\_spiketrend\_softmax & sea\_mstcn\_sep\_spiketrend & sea\_softmax\_conjunctive & ours & 0.878 & 1.43 & 0.665 & 0.488 & 0 & -- & 67,021 & 1.53 \\
55 & seanet\_bottleneck\_adaptive & sea\_mstcn\_sep\_bottleneck & sea\_adaptive\_classwise & ours & 0.868 & 1.97 & 0.680 & 0.559 & 6 & 0.7772 & 41,334 & 1.43 \\
56 & seanet\_slim\_dualstream & sea\_mstcn\_sep & sea\_dualstream\_conjunctive & ours & 0.860 & 2.40 & 0.745 & 0.498 & 0 & -- & 58,101 & 1.49 \\
57 & seanet\_gated\_mean\_adaptive & sea\_mstcn\_sep\_gated & sea\_adaptive\_classwise & ours & 0.860 & 2.24 & 0.731 & 0.549 & 0 & -- & 61,750 & 1.50 \\
58 & seanet\_inputgate\_attnmax & sea\_mstcn\_sep\_inputgate & sea\_attention\_max & ours & 0.858 & 2.15 & 0.700 & 0.625 & 0 & -- & 58,102 & 1.49 \\
59 & seanet\_spiketrend\_attnmax & sea\_mstcn\_sep\_spiketrend & sea\_attention\_max & ours & 0.854 & 1.65 & 0.644 & 0.637 & 0 & -- & 67,030 & 1.53 \\
60 & seanet\_gated\_last\_attnmax & sea\_mstcn\_sep\_gated & sea\_attention\_max & ours & 0.854 & 1.78 & 0.668 & 0.598 & 0 & -- & 61,750 & 1.50 \\
61 & seanet\_bottleneck\_shallow & sea\_mstcn\_sep\_bottleneck & sea\_topk\_conjunctive & ours & 0.846 & 2.82 & 0.624 & 0.579 & 0 & -- & 21,548 & 1.33 \\
62 & seanet\_gated\_mean\_attnmax & sea\_mstcn\_sep\_gated & sea\_attention\_max & ours & 0.832 & 1.67 & 0.704 & 0.665 & 0 & -- & 61,750 & 1.50 \\
63 & seanet\_gated\_pyramid & sea\_multiscale\_pyramid & sea\_topk\_conjunctive & ours & 0.832 & 0.98 & 0.591 & 0.602 & 0 & -- & 62,781 & 1.52 \\
64 & seanet\_bottleneck\_dualstream & sea\_mstcn\_sep\_bottleneck & sea\_dualstream\_conjunctive & ours & 0.824 & 2.04 & 0.652 & 0.638 & 0 & -- & 41,845 & 1.43 \\
65 & seanet\_bottleneck\_attnmax & sea\_mstcn\_sep\_bottleneck & sea\_attention\_max & ours & 0.816 & 1.74 & 0.610 & 0.732 & 0 & -- & 41,334 & 1.43 \\
66 & seanet\_bottleneck\_pyramid & sea\_multiscale\_pyramid & sea\_topk\_conjunctive & ours & 0.786 & 1.03 & 0.569 & 0.759 & 0 & -- & 42,365 & 1.44 \\
67 & resnet & mil\_resnet & mil\_conjunctive & millet & 0.772 & 2.95 & 0.554 & 0.824 & 85 & 0.8146 & 506,331 & 4.42 \\
68 & fcn & mil\_fcn & mil\_conjunctive & millet & 0.742 & 3.83 & 0.533 & 0.910 & 85 & 0.8141 & 267,035 & 3.47 \\
69 & seanet\_gated\_last\_dualstream & sea\_mstcn\_sep\_gated & sea\_dualstream\_conjunctive & ours & 0.736 & 1.95 & 0.637 & 0.804 & 0 & -- & 62,261 & 1.50 \\
70 & seanet\_slim\_attnmax & sea\_mstcn\_sep & sea\_attention\_max & ours & 0.714 & 1.63 & 0.532 & 0.855 & 0 & -- & 57,590 & 1.49 \\
71 & seanet\_bottleneck\_mschan\_pyramid & sea\_multiscale\_pyramid & sea\_topk\_conjunctive & ours & 0.696 & 2.17 & 0.653 & 0.992 & 0 & -- & 28,441 & 1.37 \\
72 & seanet\_spiketrend\_dualstream & sea\_mstcn\_sep\_spiketrend & sea\_dualstream\_conjunctive & ours & 0.556 & 1.64 & 0.510 & 1.260 & 0 & -- & 67,541 & 1.53 \\
\bottomrule
\end{tabular}
\end{table}

\note{I have NOT replaced this table with the \texttt{\textbackslash input} above, and
left it commented out on purpose. \texttt{table\_appendix\_full\_leaderboard.tex}'s own
caption says ``all 57 evaluated models'', but
\texttt{results/SEA\_NET/leaderboard.csv} -- the file it was generated from -- now has
66 rows: 9 more models finished training after this table was last built (the same
staleness I flagged for \texttt{table\_ranking\_accuracy} in \Cref{sec:res:ucr}).
Run \texttt{python main.py paper} once more, check the caption says 66, and only then
uncomment the \texttt{\textbackslash input} line above and delete the placeholder table
below it -- that is genuinely less error-prone than me hand-typing all 66 rows here.}

\begin{table}[h]
\caption{All 66 model variants, ranked by \webtraffic{} accuracy. \todo{Fill from
\texttt{results/SEA\_NET/leaderboard.csv}, or replace this whole table by the
generated one, \texttt{table\_appendix\_full\_leaderboard} (see the README for the
\texttt{\textbackslash input} path) -- once it has been regenerated, see the note
above.}}
\label{tab:full_leaderboard}
\begin{center}
\scriptsize
\begin{tabular}{clcccccr}
\toprule
\textbf{\#} & \textbf{Model (encoder + head)} & \textbf{Acc} & \textbf{Loss} &
\textbf{AOPCR} & \textbf{NDCG} & \textbf{Params} & \textbf{Size (MB)} \\
\midrule
1 & \texttt{seanet\_gated\_mean\_topk} (sea\_mstcn\_sep\_gated + sea\_topk\_conjunctive) & 0.954 & 0.255 & 2.293 & 0.772 & 61{,}740 & 1.50 \\
2 & \texttt{seanet\_conjunctive} (sea\_mstcn\_sep + mil\_conjunctive) & 0.954 & 0.262 & 1.502 & 0.698 & 269{,}083 & 3.54 \\
$\vdots$ & & & & & & & \\
\bottomrule
\end{tabular}
\end{center}
\end{table}
% Rows 1-2 are the two rank-1/2 entries of results/SEA_NET/leaderboard.csv (real numbers, checked
% against the CSV directly) -- kept as a concrete example of the format rather than left as \tbd,
% since a "$\vdots$ with no real row at all" table is not useful to look at. Fill the rest (or
% switch to the \input once regenerated, per the note above).

% ---------------------------------------------------------------------
\subsection{Extended ablation studies}
\label{app:ablations}

\begin{guide}
The main text keeps one ablation. Put the rest here, each as a small table or figure
with two sentences: what was varied, and what it showed.
\textbf{Candidates:} encoder depth and width; the dilation cap; the top-$k$ fraction
$\kappa$; the initial sharpness $\beta_0$; positional encoding on/off; the attention
focus penalty on/off; the attention width $\dattn$.
\end{guide}

\paragraph{A. Effect of the pooling head, encoder fixed.}
\todo{One sentence + table or figure.}

{\Cref{tab:ablation_pooling} gives the full version of the by-head marginal means that
\Cref{sec:res:ablation} only summarises in prose: averaged over every encoder it was
paired with, the pooling head alone moves \webtraffic{} accuracy by almost 18 points,
from Dual-stream Conjunctive (0.744) to Class-wise Conjunctive (0.921).}

\begin{table}[h]
\caption{\webtraffic{} accuracy, AOPCR and NDCG@$n$ averaged over every encoder each
pooling head was paired with. $n$ = number of encoder pairings averaged; the two
$n=1$ rows are the original \millet{} heads, each only ever paired with
InceptionTime, and are shown for reference rather than as part of the ranking.}
\label{tab:ablation_pooling}
\begin{center}
\small
\begin{tabular}{lcccr}
\toprule
\textbf{Pooling head} & \textbf{Mean accuracy} & \textbf{Mean AOPCR} & \textbf{Mean NDCG@$n$} & \textbf{$n$} \\
\midrule
Class-wise conjunctive (\Cref{sec:pool:classwise})   & 0.921 & 2.322 & 0.692 & 10 \\
Gated attention (\Cref{sec:pool:gated})               & 0.918 & 1.869 & 0.733 & 4 \\
Adaptive class-wise (\Cref{sec:pool:adaptive})        & 0.912 & 2.210 & 0.732 & 9 \\
Softmax conjunctive                                   & 0.907 & 1.690 & 0.712 & 9 \\
Top-$k$ conjunctive (\Cref{sec:pool:topk})            & 0.904 & 2.272 & 0.713 & 16 \\
Attention-max                                         & 0.838 & 1.758 & 0.645 & 8 \\
Dual-stream conjunctive (\Cref{sec:pool:dsmil})       & 0.744 & 2.007 & 0.636 & 4 \\
\midrule
\millet{} Additive (\citealp{early2024millet})        & 0.942 & 1.579 & 0.729 & 1 \\
\millet{} Attention (\citealp{early2024millet})       & 0.888 & 2.403 & 0.710 & 1 \\
\millet{} Conjunctive (\citealp{early2024millet})     & 0.839 & 2.712 & 0.616 & 4 \\
\bottomrule
\end{tabular}
\end{center}
\end{table}
% Source: results/SEA_NET/leaderboard.csv, grouped by the `pooling` column, mean of web_acc/
% web_aopcr/web_ndcg, n = group size (66 rows total). mil_conjunctive has n=4 rather than n=1
% because it is also the pooling head used by the re-trained ResNet/FCN/InceptionTime baselines,
% not only by MILLET's own InceptionTime encoder. Numbers match sec:res:ablation's own marginal-mean
% comment to 2-3 decimals (that comment rounds slightly differently); this table is the fuller
% version with AOPCR and NDCG added, computed directly for this appendix.

\paragraph{B. Effect of the encoder, pooling head fixed.}
\todo{One sentence + table or figure.}

{\Cref{tab:ablation_encoder} does the same for the encoder side: the spread here
(0.937 down to 0.742) is similar in size to the pooling-head spread above, which is
why \Cref{sec:res:ablation} reads the two effects as comparably important rather than
one dominating the other.}

\begin{table}[h]
\caption{\webtraffic{} accuracy, AOPCR and NDCG@$n$ averaged over every pooling head
each encoder was paired with. $n=1$ rows are the original \millet{}/baseline encoders,
each only ever paired with Conjunctive pooling.}
\label{tab:ablation_encoder}
\begin{center}
\small
\begin{tabular}{lcccr}
\toprule
\textbf{Encoder} & \textbf{Mean accuracy} & \textbf{Mean AOPCR} & \textbf{Mean NDCG@$n$} & \textbf{$n$} \\
\midrule
\texttt{sea\_multiscale\_channels}                     & 0.937 & 2.385 & 0.747 & 4 \\
\texttt{sea\_mstcn\_sep\_recon}                         & 0.924 & 2.315 & 0.719 & 5 \\
\texttt{sea\_mstcn\_sep\_inputgate} (\Cref{sec:arch:variants}) & 0.904 & 2.345 & 0.718 & 5 \\
\texttt{sea\_mstcn\_sep\_gated} (\Cref{sec:arch:variants})     & 0.902 & 2.059 & 0.711 & 19 \\
\texttt{sea\_mstcn\_sep} (\Cref{sec:arch:block})               & 0.898 & 1.933 & 0.694 & 12 \\
\texttt{sea\_mstcn\_sep\_bottleneck} (\Cref{sec:arch:variants})& 0.884 & 2.209 & 0.694 & 8 \\
\texttt{sea\_mstcn\_sep\_spiketrend}                    & 0.858 & 1.814 & 0.677 & 7 \\
\texttt{sea\_multiscale\_pyramid}                       & 0.771 & 1.394 & 0.604 & 3 \\
\midrule
InceptionTime (\millet{})                               & 0.887 & 2.569 & 0.677 & 1 \\
ResNet                                                   & 0.772 & 2.952 & 0.554 & 1 \\
FCN                                                      & 0.742 & 3.826 & 0.533 & 1 \\
\bottomrule
\end{tabular}
\end{center}
\end{table}
% Source: same leaderboard.csv, grouped by the `encoder` column. sea_mstcn_sep_gated has by far the
% largest n (19) because it is the encoder used in seanet_gated_mean_topk, the overall best model
% (Cref{fig:pareto}), and was swept against more pooling heads as a result -- worth remembering when
% comparing its mean against the other rows' smaller n.

\paragraph{C. Effect of the top-$k$ fraction $\kappa$.}
\todo{One sentence + small table: $\kappa \in \{0.05, 0.1, 0.25, 0.5, 1.0\}$ against
accuracy and AOPCR. This is the ablation that directly tests the claim of
\Cref{prop:topk}.}

\note{I did not run this sweep. \texttt{configs/models/*/*.yaml} fixes $\kappa=0.1$ for
every Top-$k$ pooling config I trained (\Cref{sec:pool:topk}); I found no record of a
$\kappa$ sweep anywhere in \texttt{results/} or in any model's \texttt{records} block.
This is a real gap: \Cref{prop:topk} is currently supported only by Top-$k$'s good
mean rank against the other heads (\Cref{tab:ablation_pooling}), not by a controlled
sweep showing accuracy/AOPCR change smoothly as $\kappa$ moves away from 0.1. If I had
more compute time left, this is the ablation I would run first -- it directly tests a
named property rather than a general comparison.}

\paragraph{D. Effect of the attention focus penalty.}
\todo{One sentence + table: with and without the penalty of \Cref{eq:loss}.}

\note{I did not run this as a controlled on/off ablation either. Every model in the
sweep was trained with the same fixed \texttt{lambda\_entropy=0.01}
(\texttt{configs/models/sv2/seanet.yaml} and every other model config); I never
trained a matched pair with \texttt{lambda\_entropy=0} to isolate its effect. The
closest evidence I have is indirect: Optuna's own search over \texttt{lambda\_entropy}
(paragraph ``Hyperparameter search'' below, range $[0, 0.05]$) picked a value close to
the default (0.0076 vs.\ 0.01) rather than 0, which is weak evidence the penalty helps
somewhat, but it is not the controlled comparison this paragraph asks for.}

% ---------------------------------------------------------------------
\subsection{Per-dataset results}
\label{app:per_dataset}

\todo{One sentence: accuracy per \ucr{} dataset for my best model and for the
re-trained baseline, so a reader can check any individual dataset.}

{\Cref{tab:per_dataset} lets a reader look up any single \ucr{} dataset behind the
mean accuracy and W/T/L record of \Cref{tab:ucr_accuracy}, rather than taking the
aggregate on faith.}

% \input{../../results/paper_figures/tables/table_appendix_per_dataset}
\note{Left commented out for the same reason as the full leaderboard above:
\texttt{table\_appendix\_per\_dataset.tex} was generated on the same stale pass
(2026-08-05) as \texttt{table\_ranking\_accuracy} and
\texttt{table\_appendix\_full\_leaderboard} -- before the \millet{} baseline's
extra-seed reruns finished (\Cref{sec:results}). Regenerate with
\texttt{python main.py paper} before uncommenting. A second reason I have not filled
the full table by hand: \Cref{tab:ucr_accuracy}'s mean is over the specific 85 datasets
\millet{} published (\texttt{main.py}'s \texttt{sweep\_order}, which I did not have
open in this session), and I only have the join over the full 128-dataset \ucr{}
archive readily available -- the two sets overlap heavily but are not identical, so
hand-picking rows risks quietly mixing them. The example rows below are drawn from the
full 128-dataset join and labelled as such, rather than presented as the 85-dataset
figure \Cref{tab:ucr_accuracy} uses.}

\begin{table}[h]
\caption{Per-dataset accuracy over the full 128-dataset \ucr{} archive (seed~$=0$;
note this is a larger set than the 85 datasets \Cref{tab:ucr_accuracy} averages over,
see the note above), a sample spanning the full range shown in \Cref{fig:acc_diff}
below. \todo{Fill the rest, ideally restricted to the 85-dataset set to match
\Cref{tab:ucr_accuracy} exactly, or \texttt{\textbackslash input} the generated table
once regenerated. Use a two-column layout so the table fits on one page.}}
\label{tab:per_dataset}
\begin{center}
\scriptsize
\begin{tabular}{lccc}
\toprule
\textbf{Dataset} & \textbf{Mine} & \textbf{\millet{} (re-trained)} & \textbf{Diff} \\
\midrule
EthanolLevel            & 0.778 & 0.280 & $+0.498$ \\
Beef                     & 0.933 & 0.733 & $+0.200$ \\
Herring                  & 0.703 & 0.531 & $+0.172$ \\
$\vdots$ & & & \\
PigCVP                   & 0.582 & 0.764 & $-0.183$ \\
SonyAIBORobotSurface1    & 0.735 & 0.957 & $-0.221$ \\
Mallat                   & 0.728 & 0.952 & $-0.224$ \\
\bottomrule
\end{tabular}
\end{center}
\end{table}
% Source: seanet_bottleneck_topk.../results.csv and millet__mil_inceptiontime__mil_conjunctive/
% results.csv, both seed 0, joined on dataset name over the full UCR archive (n=124 datasets have
% a match on both sides out of 128 -- a few are missing a seed-0 row on one side). These six rows
% are the three biggest wins and three biggest losses for seanet_bottleneck_topk against the
% re-trained MILLET baseline over that 128-set join; the same six datasets anchor fig:acc_diff
% below. I did NOT verify all six fall inside the specific 85-dataset MILLET published (most likely
% do, since 85 of the 128 are that set), which is exactly why the caption above flags the set
% difference instead of silently presenting these as tab:ucr_accuracy's 85.

% ---------------------------------------------------------------------
\subsection{Per-seed results}
\label{app:seeds}

\begin{guide}
The main text reports mean $\pm$ std. Here show the individual seeds, so the reader
can see whether the spread comes from one bad run or is uniform. Two or three
sentences are enough.
\end{guide}

{The three \webtraffic{} models with more than one seed are shown individually in
\Cref{tab:seeds}: \texttt{seanet\_bottleneck\_topk} is the tightest of the three
(0.938--0.966), the re-trained \millet{} baseline is almost as tight (0.876--0.894),
and \texttt{seanet\_inputgate\_adaptive} is the clear outlier -- one seed (0.938) is
close to \texttt{seanet\_bottleneck\_topk}'s range, but the other two (0.880, 0.898)
sit well below it, so its reported std is driven by a genuine spread across seeds, not
by one implausible run. This matches the point already made in
\Cref{sec:res:web}: the spread is not uniform across models, and the model I would
trust most from a single seed is \texttt{seanet\_bottleneck\_topk}.}

\begin{table}[h]
\caption{Individual seed results (accuracy) on \webtraffic{}.}
\label{tab:seeds}
\begin{center}
\small
\begin{tabular}{lcccc}
\toprule
\textbf{Model} & \textbf{Seed 0} & \textbf{Seed 1} & \textbf{Seed 2} &
\textbf{Mean $\pm$ std} \\
\midrule
\texttt{seanet\_bottleneck\_topk}    & 0.938 & 0.938 & 0.966 & $0.947\pm0.016$ \\
\texttt{seanet\_inputgate\_adaptive} & 0.938 & 0.880 & 0.898 & $0.905\pm0.030$ \\
\millet{} re-trained by me            & 0.894 & 0.876 & 0.892 & $0.887\pm0.010$ \\
\bottomrule
\end{tabular}
\end{center}
\end{table}
% Source: WebTraffic rows of each model's results.csv, one per seed (0/1/2). Means/stds match
% Table 7 (tab:web_main) in the main text exactly -- this table is just the per-seed breakdown
% behind that table's mean +- std columns, as the guide above asks for.

% ---------------------------------------------------------------------
\subsection{Extra experiments}
\label{app:extra}

\begin{guide}
Experiments I ran that are interesting but off the main line. One short paragraph
each; drop any that does not add something.
\end{guide}

\paragraph{Ensemble voting.}
\todo{Two or three sentences: combining the predictions of several of my variants by
soft and hard voting, and how much accuracy that buys. Say whether the gain is worth
the extra cost -- an ensemble is no longer a ``small'' model.}

{Ensembling my two proposed variants buys a real but modest gain over either alone:
soft voting reaches 0.8335 mean accuracy on the 85-dataset \ucr{} archive, against
0.8237 for \texttt{seanet\_bottleneck\_topk} and 0.8253 for
\texttt{seanet\_inputgate\_adaptive} individually (\Cref{tab:ensemble}) -- about one
accuracy point, similar in size to the gap between my re-trained \millet{} baseline
and my best single \seanet{} variant in \Cref{tab:ucr_accuracy}. Whether that point is
worth it depends entirely on the deployment target: for the on-device use case that
motivated O6 (\Cref{tab:objectives}), running both models and combining their outputs
roughly doubles the parameter count and inference cost of the already-small single
model, which works against the whole point of \Cref{sec:arch:goals}'s Goal G1; for a
setting without that constraint, a one-point gain for two small models is a reasonable
trade. Soft voting beats hard voting here (0.8335 vs.\ 0.8311), consistent with soft
voting keeping more information (the full probability, not just the arg-max) from each
member.}

\begin{table}[h]
\caption{Ensemble of my two proposed variants against each individually, mean accuracy
over the 85-dataset \ucr{} archive.}
\label{tab:ensemble}
\begin{center}
\small
\begin{tabular}{lcc}
\toprule
\textbf{Configuration} & \textbf{Mean accuracy} & \textbf{\#members} \\
\midrule
\texttt{seanet\_bottleneck\_topk} alone     & 0.8237 & 1 \\
\texttt{seanet\_inputgate\_adaptive} alone  & 0.8253 & 1 \\
Soft vote                                    & 0.8335 & 4 \\
Hard vote                                    & 0.8311 & 4 \\
\bottomrule
\end{tabular}
\end{center}
\end{table}
% Source: results/SEA_NET/ensemble_vote.csv, column means over all 85 rows (acc__seanet_
% bottleneck_topk, acc__seanet_inputgate_adaptive, acc__soft_vote, acc__hard_vote). The CSV's own
% n_members column says 4 for every row, but only 2 of the 4 member models are individually named
% in its columns -- I do not know which two additional models make up the other half of the
% ensemble, so "4" in the table above is the CSV's own stated count, not something I can break down
% further. Worth flagging rather than silently presenting a 2-model story as a 4-model one.

\paragraph{Hyperparameter search.}
\todo{Two or three sentences: what I searched with Optuna \citep{akiba2019optuna},
over what range, and whether the found setting beat the default. If it did not, say
so -- that is a useful result about how sensitive the model is.}

{For the base \seanet{} recipe (\texttt{sea\_mstcn\_sep} encoder with \millet{}'s
Additive pooling head, \Cref{sec:setup}), I ran a 30-trial Optuna search
\citep{akiba2019optuna} with a TPE sampler over nine hyperparameters at once --
learning rate, weight decay, label smoothing and the attention entropy penalty from
the training recipe, plus encoder width, depth, dropout and dilation cap, plus the
pooling head's attention width -- each trial trained on \webtraffic{} and Optuna
minimised validation loss. It did not beat the hand-set default: the best trial
reached 0.906 test accuracy, against 0.942 for the default recipe I had already been
using (\Cref{tab:cost} and \Cref{sec:setup} use the default throughout this report).
I read this as a useful negative result rather than a failed search: it suggests the
default recipe -- arrived at by hand, one change at a time, while building the
architecture in \Cref{sec:architecture} -- was already close to a local optimum for
this dataset, at least within the ranges I gave Optuna to search, so an automated
search over a fixed range was not going to find something meaningfully better.}
% Source: configs/models/sv2/seanet.yaml, records.optuna_best.metrics.test_acc = 0.906 (30 trials,
% search_space block lists all 9 tuned parameters) vs records.default.metrics.test_acc = 0.942.
% Both trained on WebTraffic per the recorded metadata (dataset: WebTraffic in both records blocks).

\paragraph{Approaches I tried and dropped.}
\todo{Two or three sentences on a direction that did not make it into the final model
(for example an earlier sliding-window formulation, or a Transformer encoder), and the
reason I dropped it.}

{One direction I scoped but did not carry through to a trained model is a Transformer
encoder (\texttt{configs/models/transformer.yaml}: \texttt{sea\_transformer}, 128-d
model, 8 heads, 4 layers, reusing \millet{}'s Additive pooling head). I want to be
precise about what this is: it is a planned-but-deferred direction, not a
trained-and-discarded one -- the config exists with draft hyperparameters and an
explicit \texttt{implemented: false} flag so the training pipeline refuses to run it
until the encoder class is actually written, and I never wrote that class. I set it
aside because the TCN-family encoders in \Cref{sec:architecture} already gave the
parameter-efficiency result O3 asked for (\Cref{tab:objectives}), and a Transformer
encoder -- with its quadratic attention cost in series length -- looked likely to move
in the opposite direction from Goal G1 (\Cref{sec:arch:goals}) before it could prove
whether it helped explanation quality. It remains a reasonable next step if a future
iteration prioritises accuracy over size.}
% Source: configs/models/transformer.yaml, read in full -- implemented: false, draft encoder/
% pooling/training blocks match the structure of the other model configs exactly, confirming this
% is a scoped-but-unbuilt placeholder rather than a model that was trained and rejected.

% ---------------------------------------------------------------------
\subsection{Additional figures}
\label{app:figures}

\begin{guide}
Figures that support the story but would slow the main text down. Each still needs a
caption that can be understood on its own.
\end{guide}

\begin{figure}[h]
  \centering
  \figph[0.8]{\textbf{Figure X -- Accuracy against series length.} One point per
  \ucr{} dataset: does my model lose accuracy on long series? This is the experiment
  that checks the dilation cap of \Cref{sec:arch:block}.}
  \caption{\texttt{seanet\_bottleneck\_topk} accuracy (seed~$=0$) against \ucr{}
  series length, one point per dataset (128 datasets, WebTraffic excluded).
  Correlation is $-0.34$: splitting the archive at the median length (344 steps), the
  shorter half averages 0.860 accuracy against 0.756 for the longer half.}
  \label{fig:acc_vs_length}
\end{figure}
% Diagram brief for fig:acc_vs_length: x-axis = series_length (log scale recommended, range 15-2844,
% from results/SEA_NET/data_summary.csv), y-axis = test_acc (0-1), one point per of the 128 UCR
% datasets (seed 0, from seanet_bottleneck_topk's results.csv joined to data_summary.csv on dataset
% name). Optionally mark the median-length split (344) with a vertical dashed line and annotate the
% two half-means (0.860 short / 0.756 long) directly on the plot -- that is the number the caption
% leans on. Pearson correlation computed directly: r = -0.3385, n = 128.
{The negative correlation is consistent with the encoder's capped dilation
(\texttt{max\_dilation=16} in every \seanet{} config, \Cref{sec:arch:block}): the
receptive field stops growing after a fixed number of steps, so on a series much
longer than that field, later blocks are working from an increasingly incomplete view
of the whole sequence. I read this as a real limitation of the design rather than
noise -- the effect size (10 points of mean accuracy between the short and long
halves) is larger than the seed-to-seed std reported anywhere in
\Cref{sec:res:training}.}

\begin{figure}[h]
  \centering
  \figph[0.8]{\textbf{Figure X -- Per-dataset accuracy difference.} One bar per
  \ucr{} dataset: my accuracy minus the re-trained baseline's, sorted. Shows
  \emph{where} I win and lose, not just how often.}
  \caption{\texttt{seanet\_bottleneck\_topk} minus re-trained \millet{} accuracy,
  seed~$=0$, one bar per \ucr{} dataset (124 of 128 with a result on both sides),
  sorted. Biggest win: EthanolLevel ($+0.498$). Biggest loss: Mallat ($-0.224$).}
  \label{fig:acc_diff}
\end{figure}
% Diagram brief for fig:acc_diff: horizontal or vertical bar chart, one bar per dataset (124 bars),
% sorted by acc_diff ascending, colour by sign (e.g. red = loss, green = win). Data: same join used
% for tab:per_dataset above (seanet_bottleneck_topk vs millet__mil_inceptiontime__mil_conjunctive
% results.csv, seed 0, dataset name). With 124 bars this will be dense -- consider labelling only
% the extremes (top/bottom 3-5) rather than every dataset name, and let the reader who wants the
% full list use tab:per_dataset or the (regenerated) full per-dataset table instead.
{The shape, not just the extremes, is the point of this figure: it is not one or two
datasets pulling a flat average up or down, but a genuine spread of wins and losses
across the whole archive, which is why the aggregate W/T/L record in
\Cref{tab:ucr_accuracy} (28/20/36) reads as close to balanced -- \Cref{fig:acc_diff}
shows that balance is real, not an artefact of how ties are counted.}

\begin{figure}[h]
  \centering
  \figph[0.8]{\textbf{Figure X -- More qualitative examples.} Four extra
  \webtraffic{} series with their importance maps, including at least one failure
  case where the explanation points at the wrong region.}
  \caption{\todo{Caption. Including a failure case is a strength -- say why it fails
  if I can tell.}}
  \label{fig:qualitative_extra}
\end{figure}
\note{Same limitation as \Cref{fig:qualitative} in the main text: this needs a live
model to generate importance maps for specific \webtraffic{} test series
(\texttt{python main.py interpret}, \Cref{sec:res:explanations}'s diagram brief), which
I cannot run in this session. Once produced, pick at least one \emph{incorrectly}
predicted or visibly poor-explanation series on purpose -- \Cref{sec:res:ucr} already
shows my proposed heads lose more often than they win on the wider \ucr{} archive, so
an appendix that only shows clean wins would understate that finding.}